\section{万有覆叠空间}
\label{ch13-sec80}

假定 $p\colon E\to B$ 是一个覆叠映射, $p(e_{0})=b_{0}$. 若 $E$ 是单连通的, 则称 $E$ 为 $B$ 的一个\textbf{万有覆叠空间}(universal covering space). 由于 $\pi_{1}(E,\ e_{0})$ 是平凡群, 在上一节中定义的对应下, 其覆叠空间也对应着 $\pi_{1}(B,\ b_{0})$ 的平凡子群. 于是 \cref{theorem:ch13-equivalence-conjugacy} 蕴涵着 $B$ 的任何两个万有覆叠空间是等价的. 基于这个原因, 我们常说 $B$ 的万有覆叠空间, 而不说 $B$ 的 (某)一个万有覆叠空间. 我们将看到, 并非每一个空间都有万有覆叠空间. 在此, 我们将在 $B$ 有万有覆叠空间的假定下讨论它的若干性质.

我们来证明以下两个基本引理.

\begin{lemma}[道路连通分支诱导覆叠映射]
\label{lemma:ch13-path-component-covering}
设 $B$ 是一个道路连通且局部道路连通的空间. 设 $p\colon E\to B$ 是较早定义的那种覆叠映射 (因此 $E$ 未必是道路连通的). 若 $E_{0}$ 为 $E$ 中的一条道路连通分支, 则由 $p$ 的限制定义的映射 $p_{0}\colon E_{0}\to B$ 是一个覆叠映射.
\end{lemma}

\begin{proof}
我们先证 $p_{0}$ 是一个满射. 由于空间 $E$ 局部同胚于 $B$, 因而 $E$ 也是局部道路连通的. 于是 $E_{0}$ 为 $E$ 中的开集. 由此可见 $p(E_{0})$ 为 $B$ 中的开集. 我们来证明 $p(E_{0})$ 也是 $B$ 中的闭集, 从而 $p(E_{0})=B$.

设 $x$ 为 $B$ 中属于 $p(E_{0})$ 闭包的一个点. 设 $U$ 是点 $x$ 的被 $p$ 均衡地覆盖的一个道路连通邻域. 由于 $U$ 包含着 $p(E_{0})$ 中的点, 那么必有 $p^{-1}(U)$ 的某一片 $V_{\alpha}$ 与 $E_{0}$ 有非空交. 由于 $V_{\alpha}$ 同胚于 $U$, 从而道路连通, 因此它必含于 $E_{0}$ 之中. 于是 $p(V_{\alpha})=U$ 包含于 $p(E_{0})$, 特别地, 有 $x\in p(E_{0})$.

现在证明 $p_{0}\colon E_{0}\to B$ 是一个覆叠映射. 给定 $x\in B$, 选取 $x$ 的邻域 $U$ 如前. 如果 $V_{\alpha}$ 是 $p^{-1}(U)$ 的一个片, 则 $V_{\alpha}$ 道路连通. 若它与 $E_{0}$ 有交, 则它包含于 $E_{0}$. 因此 $p_{0}^{-1}(U)$ 等于 $p^{-1}(U)$ 中那些与 $E_{0}$ 有交的片 $V_{\alpha}$ 的并, 每一个这样的片为 $E_{0}$ 中的开集并且被 $p_{0}$ 同胚地映到 $U$ 上. 于是 $U$ 被 $p_{0}$ 均衡地覆盖.
\end{proof}

\begin{lemma}[覆叠映射的合成]
\label{lemma:ch13-composition-covering}
设 $p$, $q$, $r$ 都是连续映射, $p=r\circ q$. 假定 $X$, $Y$, $Z$ 都是道路连通和局部道路连通空间.
\begin{enumerate}[(a)]
\item 若 $p$ 和 $r$ 都是覆叠映射, 则 $q$ 也是覆叠映射.
\item 若 $p$ 和 $q$ 都是覆叠映射, 则 $r$ 也是覆叠映射.
\end{enumerate}
\end{lemma}

\begin{proof}
(a) 假设 $p$ 和 $r$ 都是覆叠映射. 我们先证明 $q$ 是一个满射. 给定 $y\in Y$, 选取 $Y$ 中从 $y_{0}$ 到 $y$ 的一条道路 $\widetilde{\alpha}$, 则 $\alpha=r\circ\widetilde{\alpha}$ 为 $Z$ 中以 $z_{0}$ 为起点的一条道路, 设 $X$ 中以 $x_{0}$ 为起点的道路 $\widetilde{\widetilde{\alpha}}$ 为 $\alpha$ 的提升, 则 $Y$ 中以 $y_{0}$ 为起点的道路 $q\circ\widetilde{\widetilde{\alpha}}$ 为 $\alpha$ 的提升. 由道路提升的唯一性, $\widetilde{\alpha}=q\circ\widetilde{\widetilde{\alpha}}$. 因此 $q$ 将 $\widetilde{\widetilde{\alpha}}$ 的终点映到 $\widetilde{\alpha}$ 的终点 $y$. 这证明 $q$ 是满射.

给定 $y\in Y$, 我们来选取点 $y$ 的一个邻域使其被 $q$ 均衡地覆盖. 设 $z=r(y)$. 由于 $p$ 和 $r$ 都是覆叠映射, 我们可以取到点 $z$ 的一个被 $p$ 和 $r$ 均衡地覆盖的道路连通邻域 $U$. 设 $V$ 为 $r^{-1}(U)$ 中含有点 $y$ 的一片. 我们来证明 $V$ 被 $q$ 均衡地覆盖. 设 $\{U_{\alpha}\}$ 为 $p^{-1}(U)$ 中的所有片构成的族. 这时 $q$ 将每一个 $U_{\alpha}$ 映入 $r^{-1}(U)$ 中. 因为 $U_{\alpha}$ 是连通的, $q$ 必定将其映入 $r^{-1}(U)$ 的某一片中. 因此, $q^{-1}(V)$ 是一些片 $U_{\alpha}$ 的并, $U_{\alpha}$ 被 $q$ 映入 $V$ 中. 易见, 每一个这样的 $U_{\alpha}$ 都被 $q$ 同胚地映射到 $V$ 上. 设 $p_{0}$, $q_{0}$, $r_{0}$ 分别为 $p$, $q$, $r$ 的限制. 由于 $p_{0}$ 和 $r_{0}$ 为同胚, 所以 $q_{0}=r_{0}^{-1}\circ p_{0}$ 也是同胚.

(b) 我们将只在习题中用到这个结论. 设 $p$ 和 $q$ 都是覆叠映射. 因为 $p=r\circ q$ 和 $p$ 都满射, 所以 $r$ 也是满射.

给定 $z\in Z$, 设 $U$ 为点 $z$ 被 $p$ 均衡地覆盖的邻域. 我们证明 $U$ 也被 $r$ 均衡地覆盖. 设 $\{V_{\beta}\}$ 为 $r^{-1}(U)$ 的全体道路分支所构成的族, 它们是 $Y$ 中两两无交的开集. 我们来证明: 对于每一个 $\beta$, 映射 $r$ 将 $V_{\beta}$ 同胚地映到 $U$ 上.

设 $\{U_{\alpha}\}$ 为由 $p^{-1}(U)$ 的所有片构成的族, 这些片是两两无交的道路连通的开集, 因此它们便是 $p^{-1}(U)$ 的所有道路分支. 由于 $U_{\alpha}$ 是连通的, 所以它在映射 $q$ 下的像必定包含于某一个 $V_{\beta}$, 从而 $q$ 将每一个 $U_{\alpha}$ 映入集合 $r^{-1}(U)$. 因此 $q^{-1}(V_{\beta})$ 是族 $\{U_{\alpha}\}$ 的某一子族之并. 根据定理 53.2 和 \cref{lemma:ch13-path-component-covering} 可见: 若 $U_{\alpha_{0}}$ 为 $q^{-1}(V_{\beta})$ 的一条道路分支, 则 $q$ 的限制 $q_{0}\colon U_{\alpha_{0}}\to V_{\beta}$ 是一个覆叠映射. 特别地, $q_{0}$ 是一个满射. 因此 $q_{0}$ 是一个同胚, 即开的连续单射. 考虑分别由 $p$, $q$, $r$ 的限制得到的映射. 由于 $p_{0}$, $q_{0}$ 为同胚, 所以 $r_{0}$ 也是同胚.
\end{proof}

\begin{theorem}[万有覆叠空间的万有性]
\label{theorem:ch13-universal-covers-all}
设 $p\colon E\to B$ 是一个覆叠映射, 其中 $E$ 是单连通的, 那么对于任何一个覆叠映射 $r\colon Y\to B$, 存在一个覆叠映射 $q\colon E\to Y$ 使得 $r\circ q=p$.
\end{theorem}

\begin{proof}
设 $b_{0}\in B$. 选取 $e_{0}$ 和 $y_{0}$ 使得 $p(e_{0})=b_{0}$, $r(y_{0})=b_{0}$. 我们应用 \cref{lemma:ch13-general-lifting} 构造 $q$. 映射 $r$ 是一个覆叠映射, 并且由于 $E$ 是单连通的, 显然有
\[
p_{*}(\pi_{1}(E,e_{0}))\subset r_{*}(\pi_{1}(Y,y_{0})).
\]
因此, 存在一个映射 $q\colon E\to Y$ 使得 $r\circ q=p$, $q(e_{0})=y_{0}$. 由前一个引理可见, $q$ 是一个覆叠映射.
\end{proof}

我们给出一个空间作为例子说明一个空间可能没有覆叠空间. 为此需要以下引理.

\begin{lemma}[单连通覆叠空间的局部平凡性]
\label{lemma:ch13-simply-connected-local-trivial}
设 $p\colon E\to B$ 是一个覆叠映射, $p(e_{0})=b_{0}$. 若 $E$ 是单连通的, 则 $b_{0}$ 有一个邻域 $U$ 使得内射 $i\colon U\to B$ 导出平凡同态
\[
i_{*}\colon\pi_{1}(U,b_{0})\longrightarrow\pi_{1}(B,b_{0}).
\]
\end{lemma}

\begin{proof}
设 $U$ 是点 $b_{0}$ 的一个被 $p$ 均衡地覆盖的邻域, 将 $p^{-1}(U)$ 分解为片, 设 $U_{\alpha}$ 含有点 $e_{0}$ 的那一片, 设 $f$ 为 $U$ 中以 $b_{0}$ 为基点的一条回路. 由于 $p$ 决定了 $U_{\alpha}$ 与 $U$ 之间的一个同胚, 回路 $f$ 可提升为 $U_{\alpha}$ 中以 $e_{0}$ 为基点的一条回路 $\widetilde{f}$. 因为 $E$ 是单连通的, 所以 $E$ 中存在 $\widetilde{f}$ 与一条常值回路之间的一个道路同伦 $\widetilde{F}$, 于是 $p\circ\widetilde{F}$ 是 $B$ 中 $f$ 与一条常值回路之间的道路同伦.
\end{proof}

\begin{example}[无限耳环没有万有覆叠空间]
\label{example:ch13-infinite-earring}
设 $X$ 是平面中我们熟悉的一个"无限耳环". 若 $C_{n}$ 是平面上以 $(1/n,\ 0)$ 为圆心、以 $1/n$ 为半径的圆周, 则 $X$ 为所有 $C_{n}$ 的并. 设 $b_{0}$ 为原点, 如果 $U$ 是点 $b_{0}$ 在 $X$ 中的一个邻域, 我们证明由内射 $i\colon U\to X$ 所诱导的基本群之间的同态是非平凡的.

对于任何一个 $n$, 存在一个收缩 $r\colon X\to C_{n}$, 使得对于每一个 $i\neq n$, $r$ 将 $C_{i}$ 映为点 $b_{0}$. 选取充分大的 $n$ 使得 $C_{n}$ 包含于 $U$. 则下面由内射所诱导的同态的图表中 $j_{*}$ 是一个单射:
\[
\pi_{1}(C_{n},b_{0})\xrightarrow{j_{*}}\pi_{1}(X,b_{0}),
\]
经由 $\pi_{1}(U,b_{0})$. 因此, $i_{*}$ 不可能是平凡的.

从以上讨论可见, 尽管 $X$ 是道路连通且局部道路连通的, 但它却没有万有覆叠空间.
\end{example}

\begin{exercise}
\label{exercise:ch13-sec80-01}
设 $q\colon X\to Y$ 和 $r\colon Y\to Z$ 是两个映射, $p=r\circ q$.
\begin{enumerate}[(a)]
\item 设 $q$ 和 $r$ 是两个覆叠映射. 证明: 若 $Z$ 有万有覆叠空间, 则 $p$ 是一个覆叠映射. 参见第 53 节的习题 4.
\item 举例说明 $q$ 和 $r$ 是覆叠映射, 而 $p$ 却不是覆叠映射.
\end{enumerate}
\end{exercise}
